\documentclass[10pt,a4paper]{article}
\usepackage[T1]{fontenc}
\usepackage[margin=2cm]{geometry}
\usepackage{isabelle,isabellesym}
\usepackage{amsmath,amssymb}
\usepackage{pdfsetup}

\urlstyle{rm}
\isabellestyle{it}

\begin{document}

\title{A Lower Bound for the Maximum Number of Minimal Cuts}
\author{Carles Mar\'in}
\maketitle

\begin{abstract}
A minimal cut of a graph is an inclusion-minimal set of vertices whose removal disconnects it, and $c(n)$ is the largest number of minimal cuts of a graph on $n$ vertices. We formalize Seymour's observation, recorded with Erd\H{o}s problem 150~\cite{EP150,Erdos1988}, that $c(3m+2)\ge 3^m$: in the graph made of $m$ internally disjoint paths of length four between two vertices, choosing one inner vertex on each path gives $3^m$ distinct minimal cuts. The statement follows the variant of Erd\H{o}s problem 150 in the formal-conjectures repository.
\end{abstract}

\paragraph{Acknowledgement.} The formalization and the text of this document were prepared with the
assistance of generative AI systems, and checked by the author.

\tableofcontents

\parindent 0pt\parskip 0.5ex

\input{session}

\bibliographystyle{abbrv}
\bibliography{root}

\end{document}
